%=====================================================================
\chapter{Research Ethics and Human Subjects in IT Research}\label{ch:ethics}
%=====================================================================
\locator{5}{Part~5 --- Rigour, Reproducibility and Ethics}

\begin{outcomes}
By the end of this chapter you will be able to:
\begin{tight}
  \item \textbf{Apply} the four principles of research ethics to a computing
    study.
  \item \textbf{Obtain} informed consent that is genuinely informed.
  \item \textbf{Assess} re-identification risk in a dataset you intend to
    share.
  \item \textbf{Judge} the ethics of scraping, repository mining and
    observational data use.
  \item \textbf{Recognise} dual-use risk and decide what to publish.
\end{tight}
\end{outcomes}

\begin{prereq}
\chref{survey} and \chref{qualitative} if you collect data from people;
\chref{quasiexp} if you mine repositories; \chref{reproducibility} for the
tension between sharing data and protecting it.
\end{prereq}

\begin{keyterms}
respect for persons $\bullet$ beneficence $\bullet$ justice $\bullet$
informed consent $\bullet$ vulnerable population $\bullet$ anonymisation
$\bullet$ pseudonymisation $\bullet$ re-identification $\bullet$ data
minimisation $\bullet$ contextual integrity $\bullet$ dual use $\bullet$
responsible disclosure
\end{keyterms}

\section{Four Principles, and the Computing Cases That Strain Them}

\begin{center}
\small
\begin{tabular}{@{}L{3.2cm}L{4.4cm}L{6.6cm}@{}}
\toprule
\textbf{Principle} & \textbf{Requires} & \textbf{Where computing strains it} \\
\midrule
Respect for persons
  & Autonomy; informed consent; protection for those with diminished autonomy
  & Data about people who never interacted with you --- commit histories,
    forum posts, telemetry \\
Beneficence
  & Maximise benefit, minimise harm
  & Harms are diffuse and delayed: a scraped dataset causes no harm today and
    enables re-identification in five years \\
Justice
  & Burdens and benefits fairly distributed
  & Subjects recruited where access is easy; benefits accrue where the
    technology is deployed, often elsewhere \\
Respect for law and public interest
  & Compliance, transparency, accountability
  & Terms of service, jurisdiction, and research that is legal in one country
    and not another \\
\bottomrule
\end{tabular}
\end{center}

\begin{paperbox}[Ethics written for our kind of research]
David Dittrich and Erin Kenneally, \emph{The Menlo Report: Ethical Principles
Guiding Information and Communication Technology Research}, US Department of
Homeland Security, 2012~\cite{dittrich2012menlo}.\oa{}

\medskip
\textbf{Why it matters.} The Belmont Report was written for biomedical
research and does not anticipate studies in which the human subjects are
unaware, unreachable, and numerous. Menlo adapts it for ICT research and adds
the fourth principle above.

\medskip
\textbf{What to extract.} Its treatment of stakeholder analysis --- identify
everyone affected, not only those you interact with --- and its guidance on
research where consent is genuinely impossible.
\end{paperbox}

\begin{regbox}[The institutional frame --- PhD \S29]
The regulations adopt the International Center for Academic Integrity's six
values --- honesty, trust, fairness, respect, responsibility, courage --- and
require that departments ``develop procedures to prevent foreseeable risks to
academic and research integrity'' and ``regularly educate and train
students/faculty and admin staff about what constitutes academic or research
misconduct''.

\medskip
Ethical approval routes differ by department. Establish yours \emph{before}
collecting anything: retrospective approval is generally not available, and a
supervisor's verbal agreement is not approval.
\end{regbox}

\section{Informed Consent}

\begin{center}
\small
\begin{tabular}{@{}L{3.0cm}L{11.3cm}@{}}
\toprule
Informed  & They understand what will happen, what is collected, and what
            becomes of it. A wall of legal text is not informing \\
Voluntary & No coercion, and no pressure of dependency. A supervisor recruiting
            their own students, or a manager their own team, is a problem even
            with a form signed \\
Competent & They can consent. Note that in many jurisdictions minors cannot \\
Ongoing   & Withdrawable, at any point, without giving a reason and without
            penalty \\
\bottomrule
\end{tabular}
\end{center}

\begin{templatebox}[C10 --- Consent form and information sheet]
One page, plain language, at the reading level of the population.

\medskip
\begin{tabular}{@{}L{0.8cm}L{13.0cm}@{}}
$\square$ & What the study is about, in two sentences without jargon \\
$\square$ & What participation involves, and how long it will take \\
$\square$ & What data is collected --- including recordings, screen capture,
            telemetry \\
$\square$ & Who will see it, in what form, and for how long it is kept \\
$\square$ & Whether it will be published or deposited, and in what form
            (\chref{reproducibility}) \\
$\square$ & Risks and benefits, stated honestly. ``No risks'' is rarely true \\
$\square$ & That participation is voluntary and withdrawal is possible without
            penalty, \emph{including how to withdraw after the session} \\
$\square$ & What happens to already-collected data on withdrawal \\
$\square$ & Contact for questions, and a separate contact for complaints \\
$\square$ & Ethical approval reference \\
$\square$ & Signature or recorded verbal consent, dated \\
\end{tabular}
\end{templatebox}

\begin{alertbox}[The power relationship most computing studies get wrong]
Recruiting your own students is the single most common ethical problem in
graduate computing research, and a consent form does not solve it. A student
whose instructor asks them to participate cannot freely decline, whatever the
form says.

\medskip
Practical mitigations: have a colleague recruit and hold the identifying
information; ensure participation and its refusal are invisible to anyone
grading them; offer an equivalent-credit alternative activity; do not recruit
from a class you are currently assessing. If none is possible, recruit
elsewhere.
\end{alertbox}

\section{Anonymisation Is Harder Than It Looks}

\begin{center}
\small
\begin{tabular}{@{}L{3.4cm}L{10.9cm}@{}}
\toprule
Pseudonymisation & Direct identifiers replaced by codes, with a key held
                   separately. Reversible --- and therefore still personal data
                   under most data protection law \\
Anonymisation    & Irreversible, such that no one can re-identify by any
                   reasonable means. Much harder than it sounds, and often not
                   achieved by the measures people believe achieve it \\
\bottomrule
\end{tabular}
\end{center}

\begin{pitfall}[Removing names is not anonymisation]
Combinations of ordinary attributes identify people. In an interview study, a
quotation attributed to ``a female senior architect at a fintech company in
Lahore'' may identify one person exactly, and her colleagues will know who she
is even if you do not.

\medskip
In software repository data the problem is worse: commit patterns, coding
style, time-zone-inflected activity and the repositories a person contributes to
are all quasi-identifiers, and the data is often already public and linkable.

\medskip
Practical steps: aggregate demographic descriptors into bands; alter or omit
non-essential identifying detail in quotations, marking that you have done so;
check whether any single participant carries a distinctive theme; and ask a
participant to review the quotations attributed to them before publication.
\end{pitfall}

\section{Data You Did Not Collect From People}

Repository mining, scraping and telemetry analysis involve human subjects who
never met you. That they did not is not an exemption.

\begin{center}
\small
\begin{tabular}{@{}L{3.4cm}L{10.9cm}@{}}
\toprule
Public is not consented
  & A developer made their code public to enable collaboration, not to be
    studied. \term{Contextual integrity} asks whether your use matches the
    norms of the context the data came from \\
Terms of service
  & Check them. Some prohibit scraping outright; some prohibit
    redistribution, which affects what you can deposit \\
Aggregation over exposure
  & Report aggregates. A paper naming individual developers as low performers
    causes real harm for no research gain \\
Deletion
  & If a user deletes their content, your copy should follow where feasible.
    Say in the deposit how you handle this \\
Ethics review
  & Many boards treat public data as exempt. Exempt from review is not exempt
    from ethics, and the analysis in your thesis should be there regardless \\
\bottomrule
\end{tabular}
\end{center}

\begin{conceptbox}[A test that works for observational data]
Would the people in your dataset, if told exactly what you are doing, regard it
as reasonable?

\medskip
It is not a legal test and it does not settle hard cases, but it separates the
easy ones quickly. Counting commit frequencies across ten thousand
repositories: almost certainly yes. Analysing individual developers' working
hours to infer whether they have a second job: almost certainly no --- and the
data is equally public in both cases, which is the point.
\end{conceptbox}

\section{Data Protection in Practice}

\begin{center}
\small
\begin{tabular}{@{}L{3.4cm}L{10.9cm}@{}}
\toprule
Minimise        & Collect only what the research questions require. Every extra
                  field is a liability with no benefit \\
Separate        & Keep the identifier key apart from the data, on different
                  media, with different access \\
Encrypt         & At rest and in transit. A laptop with transcripts on it is a
                  breach waiting to happen \\
Retain, then delete & State a retention period in the consent form and honour
                  it. Note that this can conflict with the raw-data submission
                  \S14(i) requires --- resolve the conflict explicitly with
                  your supervisor and say so on the consent form \\
Transfer        & Cross-border transfer of personal data has legal constraints;
                  check before using an overseas cloud service \\
Breach          & Have a plan, and know who to notify \\
\bottomrule
\end{tabular}
\end{center}

\section{Dual Use}

Security research, and increasingly AI research, produces knowledge that
enables harm as well as defence.

\begin{center}
\small
\begin{tabular}{@{}L{3.6cm}L{10.7cm}@{}}
\toprule
Responsible disclosure
  & Notify the vendor, agree an embargo, publish after a fix is available.
    Standard practice in vulnerability research \\
Withholding detail
  & Publish the finding and the defence, withhold a working exploit. Contested
    --- it impedes verification, and \chref{reproducibility} would object \\
Staged release
  & Release capability progressively while monitoring misuse \\
Not publishing
  & Rare, and a serious decision. It should be made with your supervisor and
    the department, not alone \\
\bottomrule
\end{tabular}
\end{center}

\begin{alertbox}[Decide the disclosure route before you find anything]
A scholar who discovers a live vulnerability in a university or partner system
mid-project, with no agreed protocol, is in a difficult position with an
incentive to do nothing. Agree the route in advance --- who is told, in what
order, on what timeline --- and record it in the protocol. It costs one
paragraph and it is the difference between a handled incident and a crisis.
\end{alertbox}

\section{Ethics as Design, Not Paperwork}

\begin{conceptbox}[Where ethics changes the study]
Treated as a form to be completed, ethics review is an obstacle. Treated as
design, it improves the work --- and the improvements are concrete:

\begin{tight}
  \item Data minimisation reduces what you must protect, and usually clarifies
    what you are actually measuring.
  \item Anticipating deposit shapes the consent form so that sharing is
    permitted --- something impossible to add later, and the commonest reason
    good data cannot be published.
  \item Recruiting outside your own students improves external validity as well
    as voluntariness (\chref{validity}).
  \item An underpowered study wastes participants' time, which makes power
    analysis an ethical obligation and not only a methodological one
    (\chref{sampling}).
\end{tight}
\end{conceptbox}

\begin{traditions}{research ethics}
\tradEMP{Consent, minimisation, and the obligation not to waste participants on
a study incapable of answering its question.}
\tradDSR{Artefacts deployed in real settings can harm real users. Evaluation in
the wild carries obligations that a laboratory benchmark does not.}
\tradQUAL{The most demanding: deep engagement, identifiable accounts, and a
relationship with participants that continues after the study. Anonymisation is
hardest exactly where description must be thickest (\chref{qda}).}
\tradFORM{Few direct issues, until a result enables an attack --- at which
point dual use applies to a theorem as much as to an exploit.}
\end{traditions}

\begin{lab}[Ethics-review your own design]
\begin{tightnum}
  \item List every stakeholder affected, including people whose data you use
    without contacting them.
  \item For each, name the plausible harm and its mitigation.
  \item Draft the consent form with template C10 in plain language, then read
    it aloud to someone outside computing and fix what they did not understand.
  \item Assess re-identification: which attribute combinations in your data
    could identify an individual?
  \item Write the data management plan: collection, storage, retention,
    deletion, deposit.
  \item Resolve the conflict between your retention promise and
    \reg{PhD Regs 2024}{\S14(i)}, and record how.
\end{tightnum}
\end{lab}

\begin{checkpoint}
\begin{tightnum}
  \item Why is a signed consent form insufficient when recruiting your own
    students?
  \item Give three quasi-identifiers in software repository data.
  \item Public data requires no consent. Give a study using only public data
    that would nonetheless be unethical, and say what makes it so.
  \item Why is an underpowered study an ethical problem and not only a
    methodological one?
\end{tightnum}
\end{checkpoint}

\begin{chaptersummary}
\begin{tight}
  \item Four principles: respect for persons, beneficence, justice, and respect
    for law and public interest. The Menlo Report adapts them to research where
    subjects are unaware and numerous.
  \item Consent must be informed, voluntary, competent and ongoing. Recruiting
    your own students is not voluntary; mitigate or recruit elsewhere.
  \item Removing names is not anonymisation. Attribute combinations identify
    people, and repository data is full of quasi-identifiers.
  \item Public data is not consented data. Ask whether the subjects would
    regard your use as reasonable.
  \item Minimise, separate, encrypt, retain then delete --- and resolve the
    conflict with \S14(i)'s raw-data requirement explicitly.
  \item Agree a disclosure route before you find a vulnerability.
  \item Ethics as design improves studies. Anticipate deposit in the consent
    form, or your data can never be shared.
\end{tight}
\end{chaptersummary}

\begin{reviewq}
  \item Ethics boards frequently exempt public-data research from review. Given
    the re-identification risks in repository data, is the exemption defensible?
    Propose a workable review process that is proportionate.
  \item \reg{PhD Regs 2024}{\S14(i)} requires raw data submission, while data
    protection practice requires deletion after a retention period. Set out the
    conflict precisely and propose a policy that satisfies both.
  \item Contextual integrity asks whether a use matches the norms of the context
    data came from. Apply it to three uses of public GitHub data, ranked from
    clearly acceptable to clearly not, and say where you would draw the line.
  \item Responsible disclosure delays publication and can conflict with a
    scholar's need to publish before submission under \S15. How should a
    department handle a student in that position?
\end{reviewq}
